\section{Literature Review}

Wi-Fi Channel State Information (CSI) has emerged as a powerful medium for passive human sensing, enabling applications in human activity recognition (HAR) and individual identification. This section reviews foundational and recent contributions to CSI-based sensing, focusing on deep learning methodologies, dataset advancements, signal processing innovations, environmental considerations, and comparative model evaluations that motivate the present study.

\subsection{Foundations of Wi-Fi CSI-Based Human Sensing}

Early research demonstrated that the rich spatio-temporal patterns in CSI signals can be effectively leveraged for human activity recognition and user identification \parencite{yousefi2017survey,mosharaf2024wifiid}. These studies utilized deep learning architectures such as BiLSTMs, CNN--GRU hybrids, InceptionTime, and classical LSTMs to extract discriminative motion signatures embedded within CSI amplitude and phase dynamics \parencite{wakili2025bilstmcnngru,zhuravchak2022wifi_har}. The availability of multi-subject, multi-environment datasets further strengthened the capacity of these models to generalize across subjects and activities \parencite{alsaify2020losnlos,yang2023sensefi}.

\subsection{Deep Learning Architectures for CSI Interpretation}

Modern CSI-based HAR and identification systems increasingly rely on hybrid and hierarchical architectures. Models combining 1D CNN blocks with global pooling, temporal encoders such as BiLSTM layers, and lightweight yet expressive blocks such as MBConv1D or MobileNetV3 modules have significantly improved feature extraction performance.

These architectures are designed to capture both local spectral variations and long-range temporal dependencies inherent in CSI data. Their success stems from their ability to detect subtle patterns in wireless signal perturbations caused by human gait, gestures, or routine activities.

\subsection{Privacy-Preserving Sensing and Biometric Identification}

CSI-based sensing offers an inherently privacy-preserving alternative to camera-based biometrics. Unlike visual systems, Wi-Fi signals do not capture identifiable visual features; instead, they infer motion patterns from ambient signal distortions \parencite{ma2019wifisensing}. Studies reveal that gait signatures extracted from CSI are extremely difficult to reproduce or spoof, making them trustworthy for biometric authentication \parencite{zeng2016wiwho,zhang2021gaitsense,chen2023wigait}. 

The sensitivity of CSI to fine-grained body movements makes it a strong candidate for robust, unobtrusive human identification, particularly in secure or privacy-sensitive environments.

\subsection{Challenges in CSI Data: Noise, Environment, and Variability}

CSI measurements are inherently influenced by multipath fading, environmental dynamics, network configurations, and hardware variability. This sensitivity necessitates advanced architectures capable of filtering noise and extracting resilient features \parencite{pereira2025sokcsi}. Prior research highlights the role of robust signal processing methods—such as amplitude denoising, CSI ratio computation, and phase sanitization—in improving feature stability \parencite{yang2021csienhance,zeng2021csiratio}.

Additionally, the reliability of CSI signals is heavily dependent on router placement, environmental layout, and channel settings \parencite{zhuravchak2022wifi_har,pereira2025sokcsi}. Understanding these dependencies is essential for developing generalizable systems suitable for deployment across diverse indoor environments.

\subsection{Comparative Model Evaluations and Benchmarking Efforts}

Several studies emphasize the importance of comparative evaluations across different architectures to identify optimal model choices under varying constraints \parencite{yang2023sensefi}. Deep learning models such as DenseNet1D, EfficientNet1D-LSTM, and CSI-LSTMNet have demonstrated promising results in capturing distinctive gait signatures from CSI streams \parencite{zhang2021gaitsense}. 

However, model performance can depend significantly on network characteristics, training strategies, and hyperparameter settings. For this reason, recent benchmarking initiatives advocate open-access datasets and shared codebases to promote transparency, reproducibility, and fair comparison across research efforts \parencite{csibench2025,yang2023sensefi}.

\subsection{Critical Analysis of Existing Work}

Despite significant advancements, several challenges remain in the deployment of CSI-based biometric systems:

\paragraph{Lack of Standardization in Data Collection.}
Most prior studies employ unique data collection setups, varying in router models, antenna configurations, spacing, and channel properties. This inconsistency complicates cross-paper comparisons and limits generalization across environments.

\paragraph{Overfitting to Controlled Environments.}
A considerable portion of existing work is evaluated in highly controlled laboratory settings. While impressive accuracies are often achieved, real-world environments involve greater noise, dynamic obstacles, and unpredictable movement, which many models fail to fully address.

\paragraph{Limited Exploration of Multi-Task and Multi-Modal Learning.}
Although hybrid CNN–LSTM and attention-enhanced networks are widely explored, multi-modal fusion (e.g., CSI + inertial sensors) or multi-task frameworks (e.g., HAR + identification jointly) are rarely studied despite their potential to improve robustness.

\paragraph{Insufficient Focus on Computational Efficiency.}
Only a few studies evaluate the feasibility of real-time deployment. Lightweight architectures such as MobileNetV3\_1D\_LSTM remain underexplored despite their suitability for embedded or IoT platforms.

\paragraph{Need for Larger and More Diverse Datasets.}
Current datasets are smaller and less diverse than those used in other biometric fields (e.g., face recognition). Larger, regionally diverse, multi-floor, multi-AP datasets are essential to reflect realistic deployment scenarios.

These limitations motivate the need for comprehensive comparative studies—such as the one undertaken in this thesis—that evaluate multiple architectures under consistent settings, assess robustness to environmental changes, and consider computational efficiency alongside accuracy.

\subsection{Motivation for the Present Study}

The gaps identified in prior research underscore the necessity of a systematic comparison of lightweight and heavyweight deep learning architectures for CSI-based gait identification. This thesis addresses these challenges by evaluating four architectures—CSILSTMNet, DenseNet1D, EfficientNet1D-LSTM, and MobileNetV3\_1D\_LSTM on a custom multi-activity, multi-subject dataset. The goal is to determine which model most effectively captures individual-specific movement signatures while maintaining computational efficiency for real-time applications \parencite{yang2023sensefi,li2025wifigeneralizability}.

Furthermore, by integrating standardized preprocessing, controlled training strategies, and environment-specific analyses, this study contributes to the broader pursuit of reliable, privacy-preserving, and scalable CSI-based biometric systems.

